% Source coverage: physical PDF pp. 468--473 (printed pp. 445--450),
% beginning at the Section 23.4 heading and ending immediately before
% Section 23.5. Numbered equations: (23.4.1)--(23.4.14).
\subsection{The Cartan--Maurer Integral Invariant}
\label{sec:23.4}

In understanding the topology of various compact manifolds, it is a great
help that there is often a topologically invariant quantity that can be
written as an integral over the manifold. This will be important in our
discussion of instantons in the next section, and it has already been used
in studying Wess--Zumino--Witten terms in
\hyperref[sec:19.8]{Sections 19.8} and \hyperref[sec:22.7]{22.7}.

Consider a mapping of an arbitrary compact manifold $\mathcal S$ of odd
dimensionality $d$ with coordinates
$\theta^1,\theta^2,\ldots,\theta^d$ into a manifold $\mathcal M$ of
matrices $g(\theta^1,\theta^2,\ldots,\theta^d)$ with
$\operatorname{Det}g\ne0$. (For the applications that concern us here,
$\mathcal S$ is usually a sphere $S_d$, and the $g$ are the elements of a
Lie group $G$ in some representation.) We define a functional of
$g(\theta)$, known as the \emph{Cartan--Maurer form}:
\begin{align}
\mathcal I[g]\equiv{}&
\int \dd\theta^1\,\dd\theta^2\cdots \dd\theta^d\,
\epsilon^{i_1i_2\cdots i_d}
\notag\\
&{}\cdot\operatorname{Tr}\left\{
g^{-1}(\theta)\frac{\partial g(\theta)}{\partial\theta^{i_1}}
g^{-1}(\theta)\frac{\partial g(\theta)}{\partial\theta^{i_2}}
\cdots
g^{-1}(\theta)\frac{\partial g(\theta)}{\partial\theta^{i_d}}
\right\}.
\tag{23.4.1}\label{eq:23.4.1}
\end{align}
Here $\epsilon^{i_1i_2\cdots i_d}$ is the totally antisymmetric quantity
with $\epsilon^{12\cdots \dd}=1$. From the fact that
$\epsilon^{i_1i_2\cdots i_d}=-(-1)^d
\epsilon^{i_2\cdots i_di_1}$, we see that $\mathcal I[g]$ vanishes where
$\mathcal S$ is even-dimensional, so we will restrict ourselves here to
the case where $d$ is odd. The usefulness of this quantity arises from
its several remarkable properties.

First, this integral is independent of the coordinate system used to
parameterize the manifold $\mathcal S$. This follows rather obviously
from the fact that $\epsilon^{i_1i_2\cdots i_d}$ is a contravariant tensor
density, in the sense that
\[
\epsilon^{i_1i_2\cdots i_d}
\frac{\partial\theta'^{j_1}}{\partial\theta^{i_1}}
\frac{\partial\theta'^{j_2}}{\partial\theta^{i_2}}\cdots
\frac{\partial\theta'^{j_d}}{\partial\theta^{i_d}}
=\operatorname{Det}\left(\frac{\partial\theta'}{\partial\theta}\right)
\epsilon^{j_1j_2\cdots j_d}.
\]

Second, the integral \eqref{eq:23.4.1} is also invariant under small
deformations of the mapping $\mathcal S\mapsto\mathcal M$. Using the
properties of the trace, we see that under an infinitesimal change
$g\to g+\delta g$ of the function $g(\theta)$, the change in each factor
$g^{-1}\partial g/\partial\theta^{i}$ in \eqref{eq:23.4.1} makes the same
contribution to the change in $\mathcal I[g]$:
\begin{align*}
\delta\mathcal I[g]={}&\dd\int \dd\theta^1\,\dd\theta^2\cdots \dd\theta^d\,
\epsilon^{i_1i_2\cdots i_d}
\\
&{}\cdot\operatorname{Tr}\left\{
g^{-1}(\theta)\frac{\partial g(\theta)}{\partial\theta^{i_1}}
g^{-1}(\theta)\frac{\partial g(\theta)}{\partial\theta^{i_2}}
\cdots
\delta\left(
g^{-1}(\theta)\frac{\partial g(\theta)}{\partial\theta^{i_d}}
\right)\right\}.
\end{align*}
Now, the last factor in the trace is
\begin{align*}
\delta\left(
g^{-1}(\theta)\frac{\partial g(\theta)}{\partial\theta^{i_d}}
\right)
={}&-g^{-1}(\theta)\delta g(\theta)g^{-1}(\theta)
\frac{\partial g(\theta)}{\partial\theta^{i_d}}
+g^{-1}(\theta)\frac{\partial\delta g(\theta)}
{\partial\theta^{i_d}}
\\
={}&g^{-1}(\theta)\frac{\partial}{\partial\theta^{i_d}}
\left(\delta g(\theta)g^{-1}(\theta)\right)g(\theta).
\end{align*}
When we integrate by parts, the derivative
$\partial/\partial\theta^{i_d}$ gives no contribution when acting on the
partial derivatives $\partial g(\theta)/\partial\theta^{i_m}$ because
$\epsilon^{i_1i_2\cdots i_d}$ is antisymmetric. The remaining $d-1$
terms where $\partial/\partial\theta^{i_d}$ acts on $g^{-1}(\theta)$ are
all equal except for an alternating sign, so since there are an even
number of them they add up to zero.

Finally, let us specialize to the case where $\mathcal S$ is the sphere
$S_d$. Because $\mathcal I[g]$ is invariant under small variations of
$g(\theta)$, it can be regarded as a function $\mathcal I(c)$ only of the
homotopy class $c$ to which $g(\theta)$ belongs. The integrals
$\mathcal I(c)$ (or strictly speaking, $\exp\{\mathcal I(c)\}$) furnish a
representation of the homotopy group $\pi_d(\mathcal M)$, in the sense
that
\begin{equation}
\mathcal I(c_a\cdot c_b)=\mathcal I(c_a)+\mathcal I(c_b).
\tag{23.4.2}\label{eq:23.4.2}
\end{equation}
(If $g_a(\theta)$ and $g_b(\theta)$ are elements of the homotopy classes
$c_a$ and $c_b$ respectively, then the homotopy class $c_a\cdot c_b$
consists of mappings homotopically equivalent to
\[
g_{ab}(\theta)=
\begin{cases}
g_a(2\theta_1,\theta_2\cdots\theta_d),
&0\leq\theta_1\leq\tfrac12,\\
g_b(2\theta_1-1,\theta_2\cdots\theta_d),
&\tfrac12\leq\theta_1\leq1.
\end{cases}
\]
The part of the integral $\mathcal I[g_{ab}]$ over the hemispheres with
$0\leq\theta_1\leq1/2$ and $1/2\leq\theta_1\leq1$ can be done by changing
variables to $\theta'_1=2\theta_1$ and $\theta'_1=2\theta_1-1$,
respectively, yielding the terms $\mathcal I(c_a)$ and $\mathcal I(c_b)$
in \eqref{eq:23.4.2}.)

In particular, this tells us that the homotopy classes
$e,c,c\cdot c$, etc., and $c^{-1},c^{-1}\cdot c^{-1}$, etc., have
\begin{equation}
\mathcal I(c^n)=n\mathcal I(c).
\tag{23.4.3}\label{eq:23.4.3}
\end{equation}
If $\mathcal I(c)\ne0$ for some $c$, then these invariants are all
different, so the classes $c^n$ are all different, and therefore form a
subgroup $\mathbb{Z}$ of $\pi_d(\mathcal M)$. This goes far to explain the great
difference between the sizes of the homotopy groups for odd and even
dimensions shown in \hyperref[app:23.B]{Appendix B} of this chapter. For instance,
$\pi_1(U(1))=\mathbb{Z}$, $\pi_3(G)=\mathbb{Z}$ for all simple Lie groups $G$, and
$\pi_5(SU(n))=\mathbb{Z}$ for all $n\geq3$, while for all Lie groups $G$,
$\pi_2(G)=0$ and $\pi_4(G)$ is finite.

As a simple example where $\mathcal I(c)\ne0$, consider the homotopy
group $\pi_1(U(1))$, which is the same as the group $\pi_1(S_1)$ used as
an example at the beginning of the previous section. Any mapping
$S_1\mapsto U(1)$ may be characterized by $\nu$, the number of times that
the phase of the $U(1)$ element goes counterclockwise around $S_1$ minus
the number of times it goes around $S_1$ in the opposite direction, as
the coordinate $\theta$ goes around $S_1$, two mappings being
homotopically equivalent if and only if they have the same $\nu$. The
$\nu$th class contains the mapping
$g_\nu(\theta)=\exp(2\ii\nu\pi\theta)$ with $0\leq\theta\leq1$, for which
\[
\mathcal I[g_\nu]
=\int_0^1\dd\theta\,\exp(-2\ii\nu\pi\theta)
\frac{\dd}{\dd\theta}\exp(2\ii\nu\pi\theta)
=2\ii\nu\pi,
\]
thus verifying that $\pi_1(U(1))=\mathbb{Z}$.

As an aid to calculating $\mathcal I(g)$ in less simple cases, suppose
that we can continuously deform the manifold $\mathcal M$ into a Lie
group $H$ of dimensionality $d$. The result of performing the $H$
transformation with parameters $\theta$ followed by the $H$
transformation with parameters $\varphi$ is an $H$ transformation, say
with parameters $\theta'(\theta,\varphi)$. In terms of a matrix
representation $g(\theta)$, this reads
\[
g(\varphi)g(\theta)=g\bigl(\theta'(\theta,\varphi)\bigr).
\]
Differentiating with respect to $\theta'$ with $\varphi$ fixed and
multiplying on the left with the inverse of this equation gives
\[
\frac{\partial\theta^i}{\partial\theta'^j}
g^{-1}(\theta)\frac{\partial g(\theta)}{\partial\theta^i}
=g^{-1}(\theta')\frac{\partial g(\theta')}{\partial\theta'^j}.
\]
The integrand of $\mathcal I[g]$ at a point $\theta'$ is therefore
\begin{align*}
&\epsilon^{j_1j_2\cdots j_d}\operatorname{Tr}\left\{
g^{-1}(\theta')\frac{\partial g(\theta')}{\partial\theta'^{j_1}}
g^{-1}(\theta')\frac{\partial g(\theta')}{\partial\theta'^{j_2}}
\cdots
g^{-1}(\theta')\frac{\partial g(\theta')}{\partial\theta'^{j_d}}
\right\}
\\
&\quad=\operatorname{Det}\left(\frac{\partial\theta}
{\partial\theta'}\right)
\epsilon^{i_1i_2\cdots i_d}\operatorname{Tr}\left\{
g^{-1}(\theta)\frac{\partial g(\theta)}{\partial\theta^{i_1}}
g^{-1}(\theta)\frac{\partial g(\theta)}{\partial\theta^{i_2}}
\cdots
g^{-1}(\theta)\frac{\partial g(\theta)}{\partial\theta^{i_d}}
\right\}.
\end{align*}

Now, every Lie group $H$ has a metric $\gamma_{ij}(\theta)$ (not
necessarily unique) which is form-invariant in the sense that
\begin{equation}
\gamma_{ij}(\theta')
=\frac{\partial\theta^k}{\partial\theta'^i}
\frac{\partial\theta^\ell}{\partial\theta'^j}
\gamma_{k\ell}(\theta).
\tag{23.4.4}\label{eq:23.4.4}
\end{equation}
For instance, we can take
\begin{equation}
\gamma_{ij}(\theta)=-\frac12\operatorname{Tr}\left\{
g^{-1}(\theta)\frac{\partial g(\theta)}{\partial\theta^i}
g^{-1}(\theta)\frac{\partial g(\theta)}{\partial\theta^j}
\right\}.
\tag{23.4.5}\label{eq:23.4.5}
\end{equation}
For any choice of $\gamma_{ij}(\theta)$, the determinant of
Eq.~\eqref{eq:23.4.4} gives
\[
\operatorname{Det}\left(\frac{\partial\theta}{\partial\theta'}\right)
=\sqrt{\frac{\operatorname{Det}\gamma(\theta')}
{\operatorname{Det}\gamma(\theta)}}.
\]
By replacing the coordinates $\theta$ in Eq.~\eqref{eq:23.4.1} with
$\theta'$, we find
\begin{align}
\mathcal I[g]={}&
\epsilon^{i_1i_2\cdots i_d}\operatorname{Tr}\left\{
g^{-1}(\theta)\frac{\partial g(\theta)}{\partial\theta^{i_1}}
g^{-1}(\theta)\frac{\partial g(\theta)}{\partial\theta^{i_2}}
\cdots
g^{-1}(\theta)\frac{\partial g(\theta)}{\partial\theta^{i_d}}
\right\}
\notag\\
&{}\cdot\frac{1}{\sqrt{\operatorname{Det}\gamma(\theta)}}
\int \dd^d\theta'\sqrt{\operatorname{Det}\gamma(\theta')}.
\tag{23.4.6}\label{eq:23.4.6}
\end{align}
Since the parameters $\varphi$ of the second $H$ transformation are
arbitrary, we may regard $\theta$ and $\theta'$ as independent variables,
and evaluate the right-hand side of \eqref{eq:23.4.6} at any value of
$\theta$, say $\theta^i=0$. It will be convenient to normalize the
generators $t_i$ and coordinates $\theta^i$ so that for $\theta\to0$,
\begin{equation}
g(\theta)\longrightarrow \mathbf{1}+2\ii\theta^it_i.
\tag{23.4.7}\label{eq:23.4.7}
\end{equation}
In this case, Eq.~\eqref{eq:23.4.6} reads
\begin{equation}
\mathcal I[g]=(2\ii)^d\epsilon^{i_1i_2\cdots i_d}
\operatorname{Tr}\{t_{i_1}t_{i_2}\cdots t_{i_d}\}
\frac{1}{\sqrt{\operatorname{Det}\gamma(0)}}
\int \dd^d\theta'\sqrt{\operatorname{Det}\gamma(\theta')}.
\tag{23.4.8}\label{eq:23.4.8}
\end{equation}

We will be especially interested in the case $d=3$.
Bott~\hyperref[ch23-ref-22]{[22]} has shown that for any simple Lie group
$G$, all continuous mappings $S_3\mapsto G$ may be continuously deformed
into mappings of $S_3$ into a `standard' $SU(2)$ subgroup of $G$. (Where
$G=SU(n)$, this standard $SU(2)$ subgroup is the one that acts only on the
first two components of the defining representation of $SU(n)$. Not all
$SU(2)$ subgroups of $SU(n)$ are equivalent to this one.) As remarked in
\hyperref[sec:2.7]{Section 2.7}, in a $2\times2$ representation the general
element of $SU(2)$ may be written
\begin{equation}
g(\theta)=
\begin{pmatrix}
\theta_4+\ii\theta_3&\theta_2+\ii\theta_1\\
-\theta_2+\ii\theta_1&\theta_4-\ii\theta_3
\end{pmatrix}
=\theta_4+2\ii\boldsymbol{\theta}\cdot\mathbf t,
\tag{23.4.9}\label{eq:23.4.9}
\end{equation}
where as usual
\[
t_1=\frac12\begin{pmatrix}0&1\\1&0\end{pmatrix},
\qquad
t_2=\frac12\begin{pmatrix}0&-\ii\\\ii&0\end{pmatrix},
\qquad
t_3=\frac12\begin{pmatrix}1&0\\0&-1\end{pmatrix},
\]
and $\theta_4$ and $\boldsymbol{\theta}$ are real, with
$(\theta_4)^2=1-\boldsymbol{\theta}^2$. (Note that
Eq.~\eqref{eq:23.4.9} is consistent with the normalization convention
\eqref{eq:23.4.7}.) A straightforward calculation gives the metric
\eqref{eq:23.4.5} as
\begin{equation}
\gamma_{ij}(\theta)=\delta_{ij}
+\frac{\theta_i\theta_j}{1-\boldsymbol{\theta}^2},
\tag{23.4.10}\label{eq:23.4.10}
\end{equation}
so that
\begin{equation}
\operatorname{Det}\gamma(\theta)
=\frac{1}{1-\boldsymbol{\theta}^2}.
\tag{23.4.11}\label{eq:23.4.11}
\end{equation}
Hence Eq.~\eqref{eq:23.4.8} here reads
\[
\mathcal I[g]=-8\ii\epsilon^{ijk}\operatorname{Tr}\{t_it_jt_k\}
\int\frac{\dd^3\theta}{\sqrt{1-\boldsymbol{\theta}^2}}.
\]
Using $4t_it_j=\delta_{ij}+2\ii\epsilon^{ij\ell}t_\ell$ and
$\operatorname{Tr}\{t_\ell t_k\}=\frac12\delta_{\ell k}$, we see that
\[
8\epsilon^{ijk}\operatorname{Tr}\{t_it_jt_k\}
=2\ii\epsilon^{ijk}\epsilon^{ijk}=12\ii.
\]
Also, for the `identity' mapping $g_1$, the integral here runs twice over
the interior of the unit ball (because $\theta_4$ can be positive or
negative), and gives
\[
\int\frac{\dd^3\theta}{\sqrt{1-\boldsymbol{\theta}^2}}
=2\int_0^1\frac{4\pi r^2\,\dd r}{\sqrt{1-r^2}}=2\pi^2.
\]
For the class $c$ of mappings homotopic to $g_1$, we have then
\begin{equation}
\mathcal I(c)=24\pi^2
\tag{23.4.12}\label{eq:23.4.12}
\end{equation}
and so
\begin{equation}
\mathcal I(c^\nu)=24\pi^2\nu.
\tag{23.4.13}\label{eq:23.4.13}
\end{equation}
The integer $\nu$ is known as the \emph{winding number}. This result is
for a representation and normalization conventions for which the
standard $SU(2)$ subalgebra has generators $t_i$ with structure constants
$\epsilon_{ijk}$ and $\operatorname{Tr}(t_it_j)=\frac12\delta_{ij}$. More
generally, if $[t_i,t_j]=ig\epsilon_{ijk}t_k$ and
$\operatorname{Tr}(t_it_j)=\frac12Ng^2\delta_{ij}$, then
\begin{equation}
\mathcal I(c^\nu)=24\pi^2N\nu.
\tag{23.4.14}\label{eq:23.4.14}
\end{equation}

The results \eqref{eq:23.4.13} or \eqref{eq:23.4.14} show incidentally
that for every simple Lie group, $\pi_3(G)$ contains $\mathbb{Z}$. As listed in
\hyperref[app:23.B]{Appendix B} to this chapter, $\pi_3(G)=\mathbb{Z}$ for all simple Lie groups. Thus
the homotopy class of $g(\theta)$ for any simple Lie group is entirely
determined by its homotopy class when the group is deformed into its
standard $SU(2)$ subgroup.
